\mxnewpage\section*{Correction EV2-T1}\label{EV2-T1-CORR}
\textbf{QCM :} 1 A, 2 D, 3 D, 4 B, 5 B, 6 C, 7 B, 8 C, 9 A, 10 C, 11 D, 12 C. Chaque réponse correcte : \num{0,5} point ; total 6.
\paragraph{Vérifications des automatismes.} Priorités ; somme de fractions ; soustraction des exposants ; coefficient multiplicateur \num{1,2} ; division par \num{0,8} ; produit \num{1,25} fois \num{0,8} ; résolution affine ; accroissement vertical divisé par l’horizontal ; moyenne ; angle divisé par 360 ; complément à 1 ; substitution entre parenthèses.
\subsection*{Exercice 1 -- Fabrication et contrôle}
\begin{enumerate}
\item Premier niveau : $\num{0,3}$ et $\num{0,7}$. Après A : $\num{0,1}$ et $\num{0,9}$ ; après B : $\num{0,2}$ et $\num{0,8}$. Chaque paire somme à 1. \textit{Barème : 1 point. Résultat et justification élémentaire attendus ; demi-point possible pour démarche correcte inachevée.}
\item $P(D)=\num{0,03}+\num{0,14}=\num{0,17}$ ; $P_D(A)=\dfrac{3}{17}$. \textit{Barème : 2 points. Méthode ou justification : 1 ; calcul et conclusion : 1.}
\item Non : $P_A(D)=\num{0,1}\ne\num{0,17}$. \textit{Barème : 2 points. Méthode ou justification : 1 ; calcul et conclusion : 1.}
\end{enumerate}
\subsection*{Exercice 2 -- Minimum d’un quotient}
\begin{enumerate}
\item $f\prime =\dfrac{(x-3)(x+3)}{x^2}$. \textit{Barème : 1 point. Résultat et justification élémentaire attendus ; demi-point possible pour démarche correcte inachevée.}
\item Décroît jusqu’à 3 puis croît ; minimum 6. \textit{Barème : 2 points. Méthode ou justification : 1 ; calcul et conclusion : 1.}
\item $1-\dfrac{9}{x^2}=\dfrac{3}{4}$ donne $x^2=36$, donc x=6 seul sur le domaine ; point $\pt{6;\dfrac{15}{2}}$. \textit{Barème : 2 points. Méthode ou justification : 1 ; calcul et conclusion : 1.}
\end{enumerate}
\subsection*{Exercice 3 -- Projeté et cercle}
\begin{enumerate}
\item $\pt{1;1}$. \textit{Barème : 1 point. Résultat et justification élémentaire attendus ; demi-point possible pour démarche correcte inachevée.}
\item $\lambda =\dfrac{5+2-3}{2}=2$ ; H(3;0). \textit{Barème : 1 point. Résultat et justification élémentaire attendus ; demi-point possible pour démarche correcte inachevée.}
\item $PH^2=8$ ; $(x-5)^2+(y-2)^2=8$. \textit{Barème : 2 points. Méthode ou justification : 1 ; calcul et conclusion : 1.}
\end{enumerate}
\paragraph{Étalonnage.} Durée cible : 20–25 min pour les automatismes, 80–85 min pour les exercices, 10–15 min de relecture. Ces durées sont des estimations éditoriales à mesurer en classe. Accepter toute méthode correcte ; ne pas sanctionner deux fois une erreur numérique propagée avec cohérence. Le barème proposé ne vaut pas décision d’harmonisation de l’équipe.
